\chapter{An Architecture of Authority}
\label{ch:13}

The principle that organizes the program receives its most concrete expression in the proposal for robotic surrogates directed by human operators. In this arrangement a physical or virtual embodiment is equipped with local autonomy for bounded functions: maintaining balance, avoiding collisions, limiting applied force, preventing injury, and executing emergency stops. These functions are well suited to automation because they are fast, their success conditions are specifiable, and their failure modes can be tested against explicit criteria. Decisions concerning teaching, caregiving, emotional development, and interpersonal interaction remain with qualified human operators, who direct the embodiment and who are responsible for what it does. A single operator might supervise several embodiments, although the safe number depends on the complexity of the task, the latency of the communication link, and the degree of continuous attention required. The allocation follows the tradition of human supervisory control~\cite{sheridan1992}, and the warning of Scott~\cite{scott1998} about schemes of legibility imposed from above applies to any such allocation that is designed without those who live under it. The arrangement is notable because it specifies an architecture of authority and does not rely on an abstract promise of beneficial behavior: the allocation of control between machine and person is explicit, and so is the allocation of responsibility. The classification of automation by stages and levels, which separates the acquisition of information, its analysis, the selection of an action, and its implementation, gives the allocation its standard vocabulary~\cite{parasuraman2000}, and the ladder of citizen participation~\cite{arnstein1969} states the corresponding distinction between consulting a public and sharing power with it.

The architecture applies, with differing limits on delegation, to education, eldercare, rehabilitation, remote medical assistance, and hazardous industrial work, and the differences among these domains are themselves informative about the principle. Where the hazards are physical and the human operator's role is to select goals, as in hazardous industrial work, the supervisory ratio may be high and the embodiment may carry substantial autonomy over motion. Where the activity is relational, as in teaching and caregiving, the value of the activity resides substantially in the quality of human attention, and the ratio must be lower. In the case of infants and very young children the relevant constraint is not a ratio at all. Remote robotic childcare should not be assumed equivalent to physically present human caregiving, since the developmental and protective functions of embodied human presence are not established to be reproducible by mediated control, and the program therefore treats physical human presence as a requirement for that class of care, with teleoperated assistance permitted only as an adjunct. Safeguarding, privacy, accountability, and auditability provisions follow from the same reasoning: because the operator is the bearer of responsibility, the operator's identity, qualifications, and actions must be recorded and reviewable, and the data collected through an embodiment in a child's environment must be governed by strict limits on retention and use.

The institutional question that this architecture answers is how increasingly capable systems might change the economics of public services without making access to those services depend on the removal of human professionals. A universal and tuition-free educational system extending from early childhood through advanced study, professional retraining, and lifelong learning is a plausible objective for such a program. Capable systems could assist with curriculum development, accessibility, translation, scheduling, and instruction in scientific subjects, and the supervised surrogate model could extend the reach of qualified educators to places and populations they currently cannot serve. In each case the improvement is a multiplication of access to human judgment rather than a replacement of it, which is the sense in which the program expands human capability without transferring authority.

A further distinction separates human direction from human accountability, and the architecture is incomplete without it. A person can authorize an action without understanding it, and an institution can formally retain responsibility for systems it cannot evaluate, so that authorization by a human is necessary but not sufficient for control in any meaningful sense. Direction is meaningful only if the directing party has adequate information about what is being proposed and what it is likely to do, has real alternatives among which to choose, including the alternative of not proceeding, can refuse without incurring a penalty that makes refusal infeasible, has access to verification performed independently of the proposer, is protected by limits that are enforced technically and not merely stated in policy, and retains the ability to stop or reverse a deployment wherever that is technically possible. The program treats these six conditions as part of what it means to say that authority remains with persons, and an arrangement in which a human signature is the only remaining human contribution fails them.

\section*{Formalization}

Let the embodiment's state evolve as $x_{t+1}=f(x_t,u_t,w_t)$ on a state space $X$, where $u_t=\pi(x_t,\xi_t)$ is the action chosen by a local controller $\pi$, $w_t\in\mathcal{W}$ is a disturbance drawn from a specified modeled set, and $\xi_t\in\Xi$ is the communication condition, including loss of link and latency up to a stated bound. Let $S\subseteq X$ be a \emph{safe set} specified in advance by human decision.

\begin{definition}[Robust invariance]
A local controller $\pi$ is \emph{robustly invariant} for $S$ from an initial set $X_0\subseteq S$ over horizon $H$ if, for every $x_0\in X_0$, every disturbance sequence $(w_t)\subseteq\mathcal{W}$, and every communication sequence $(\xi_t)\subseteq\Xi$, the resulting trajectory satisfies $x_t\in S$ for $0\le t\le H$.
\end{definition}

\begin{definition}[Delegation set]
The delegation set $\mathcal{D}$ is the set of local behaviors $\pi$ such that (a) $\pi$ is robustly invariant for $S$, supported by a certificate (a proof, a verified simulation envelope, or a tested bound) that is itself subject to independent verification, and (b) $\pi$ performs no task with $\kappa(\alpha)=\mathsf{J}$. Behaviors outside $\mathcal{D}$ require operator authorization.
\end{definition}

The delegation set thus encodes the principle that what can be delegated without surrendering what ought to remain under human direction is automated: condition (a) restricts automation to behaviors whose safety can be guaranteed, and condition (b) removes the judgment tasks regardless of whether they could be performed safely.

\begin{remark}[Reversibility]
Let $\theta(\alpha)$ denote the time after which the effects of action $\alpha$ can no longer be undone, $t_r$ the operator response time, $t_x$ the time needed to execute a reversal, and $\Delta\ell$ the worst-case uncertainty in communication latency. For classes of action for which reversal forms part of the safety case, reversibility is imposed as an additional, conjunctive requirement $\theta(\alpha)>t_r+t_x+\Delta\ell$, and it never substitutes for invariance. An action whose consequences can leave $S$ under an admissible disturbance, as when an embodiment applies force to a person, is excluded from $\mathcal{D}$ whether or not the resulting injury could later be treated.
\end{remark}

\begin{definition}[Meaningful direction]
A directing party $H$ exercises \emph{meaningful direction} over a deployment if six predicates hold jointly: $\mathrm{Info}(H)$, that $H$ has access to an evaluation of the proposal and its likely effects whose error rates are bounded at a stated level; $\mathrm{Alt}(H)$, that $H$ can choose among at least two admissible options, one of which is not to proceed; $\mathrm{Ref}(H)$, that refusal is within $H$'s authority and is not penalized in a way that makes it infeasible; $\mathrm{Ver}(H)$, that the verification predicates $\varphi_j$ relevant to the decision are operated independently of the proposer; $\mathrm{Lim}(H)$, that the constraints defining $A$ are enforced technically and not only by policy; and $\mathrm{Stop}(H)$, that $H$ can halt or reverse the deployment wherever this is technically possible.
\end{definition}

Authorization by $H$, meaning that $H$ signs off, is necessary but not sufficient for meaningful direction, since it entails none of the six predicates.

For the supervisory ratio, suppose each of $n$ embodiments places an attention demand $X_i\in[0,1]$ on the operator, as a fraction of the operator's attention, with mean $m>0$ and standard deviation $s\ge0$, and with pairwise correlation $r\in[0,1]$. The operator's capacity is $1$. The sum $\Sigma_n=\sum_iX_i$ has mean $nm$ and variance $\sigma_n^2=ns^2\big(1+(n-1)r\big)$. For a tolerance $\delta\in(0,\tfrac12]$ let $k_\delta=\sqrt{(1-\delta)/\delta}$ and let $z_{1-\delta}\ge0$ be the $(1-\delta)$ quantile of the standard normal distribution.

\begin{proposition}
(a) If $nm+k_\delta\,\sigma_n\le1$ then $\Pr(\Sigma_n\le1)\ge1-\delta$, with no distributional assumption beyond the stated mean and variance. (b) Under a normal approximation to $\Sigma_n$, the constraint $\Pr(\Sigma_n\le1)\ge1-\delta$ holds approximately when $nm+z_{1-\delta}\,\sigma_n\le1$. (c) For $r=0$ and either multiplier $\lambda\in\{k_\delta,z_{1-\delta}\}$, the largest integer $n$ satisfying $nm+\lambda s\sqrt n\le1$ is $\lfloor u^2\rfloor$ with
\[
u=\frac{-\lambda s+\sqrt{\lambda^2s^2+4m}}{2m}.
\]
(d) For $r>0$ the ratio $\sigma_n/n\to s\sqrt r$ as $n\to\infty$, so the constraint is asymptotically $n\lesssim 1/(m+\lambda s\sqrt r)$, which equals $1/(m+\lambda s)$ only when $r=1$, that is, for perfectly correlated demands.
\end{proposition}

\begin{proof}
(a) Cantelli's inequality gives $\Pr(\Sigma_n-nm\ge t)\le\sigma_n^2/(\sigma_n^2+t^2)$ for $t>0$. With $t=1-nm\ge k_\delta\sigma_n$ the right side is at most $1/(1+k_\delta^2)=\delta$, and the case $\sigma_n=0$ is trivial. (b) is the corresponding statement for the normal quantile; because each $X_i$ is bounded and the approximation may be poor in the tails and for small $n$, (a) is the safer statement. (c) Substituting $u=\sqrt n$ gives $mu^2+\lambda su-1\le0$. Since $m>0$ and $\lambda s\ge0$ the quadratic has one positive root, $nm+\lambda s\sqrt n$ is increasing in $n$, and the feasible integers therefore form an initial segment ending at $\lfloor u^2\rfloor$. (d) We have $\sigma_n^2/n^2=s^2\big(1/n+(1-1/n)r\big)\to s^2r$, and dividing the constraint by $n$ gives the stated limit.
\end{proof}

The role of the correlation parameter is central. When emergencies in separate embodiments are triggered by a common cause, such as a fire alarm or a network failure, $r$ is large, the pooling gain that makes the $\sqrt n$ scaling of case (c) available disappears, and a ratio justified under independence can be unsafe under common-cause demand. For the class of infant care add the constraint $n_{\mathrm{present}}\ge1$, a requirement that a qualified human be physically present, which cannot be satisfied by any value of the supervision ratio.
